\par\Needspace{8\baselineskip}\bigskip
\pdfbookmark[3]{Darstellung, 3}{pallas-II-article3}
\begin{center}3.\end{center}
\noindent\hspace*{\parindent}Die Summe der Gleichungen 8, 9, 10, jeweils mit \(x,y,z\) multipliziert, liefert
\[
\begin{array}{rl}20)\quad&Fx+Gy+Hz=\dfrac{pr\mathrm{d}r}{\mathrm{d}t}\mathpunct{\cdot}\end{array}
\]
\par Multipliziert man diese Gleichung mit \(F\), Gleichung 12 mit \(H\) und Gleichung 13 mit \(-G\) und addiert, so ergibt sich
\[
(FF+GG+HH)x=(BH-CG)(r-p)+\frac{Fpr\mathrm{d}r}{\mathrm{d}t}\mathpunct{\raisebox{.45ex}{,}}
\]
oder, wegen \(FF+GG+HH=\left(1-\frac{p}{a}\right)mp=mpee\),
\[
\begin{array}{rl}21)\quad&x=\dfrac{BH-CG}{mpee}(r-p)+\dfrac{Fr\mathrm{d}r}{mee\mathrm{d}t}\mathpunct{\cdot}\end{array}
\]
\par Ebenso erhält man
\[
\begin{array}{rl}
22)\quad&y=\dfrac{CF-AH}{mpee}(r-p)+\dfrac{Gr\mathrm{d}r}{mee\mathrm{d}t}\\[6pt]
23)\quad&z=\dfrac{AG-BF}{mpee}(r-p)+\dfrac{Hr\mathrm{d}r}{mee\mathrm{d}t}\mathpunct{\cdot}
\end{array}
\]
\par Nun gilt
\[
\begin{gathered}
r-p=aee-ae\cos E\\[4pt]
r\mathrm{d}r=aer\sin E\mathbin{.}\mathrm{d}E=\BareRoot{am}\mathbin{.}e\sin E\mathbin{.}\mathrm{d}t;
\end{gathered}
\]
also erhält man
\[
\begin{array}{rl}
24)\quad&x=\dfrac{a(BH-CG)}{mp}-\dfrac{a(BH-CG)}{mpe}\cos E+\dfrac{\BareRoot{am}\mathbin{.}F}{me}\sin E\\[7pt]
25)\quad&y=\dfrac{a(CF-AH)}{mp}-\dfrac{a(CF-AH)}{mpe}\cos E+\dfrac{\BareRoot{am}\mathbin{.}G}{me}\sin E\\[7pt]
26)\quad&z=\dfrac{a(AG-BF)}{mp}-\dfrac{a(AG-BF)}{mpe}\cos E+\dfrac{\BareRoot{am}\mathbin{.}H}{me}\sin E.
\end{array}
\]
\par Die vier Gleichungen 19, 24, 25, 26 enthalten die vollständige Lösung der Gleichungen 1, 2, 3.
